% TEMPLATE-MILC-v3.tex - plantilla maestra UNICA de las guias MILC v3.
%
% La usan las cuatro familias de guias, cada una con su driver:
%   · Grados 6-11          -> scripts/build-guias-g11.py
%   · Bebras               -> scripts/build-guias-bebras.py
%   · Semillero            -> scripts/build-guias-semillero.py
%   · Territorio interior  -> scripts/build-guias-territorio-interior.py
%
% Antes habia tres copias casi identicas (~400 lineas iguales, 6 distintas):
% cada mejora habia que aplicarla tres veces, y en la practica no se hacia
% (el motor de recursos vivia solo en dos de ellas). Lo que varia entre
% familias esta parametrizado en 6 placeholders que cada driver llena
% (se nombran aqui SIN su sintaxis real: el builder reemplaza por texto plano,
% tambien dentro de los comentarios, y un valor multilinea rompe el preambulo):
%
%   HEADER_IZQ        encabezado izquierdo de pagina
%   HEADER_DER        encabezado derecho de pagina
%   PORTADA_ETIQUETA  rotulo sobre el numero grande (GRADO/MOMENTO/MODULO)
%   PORTADA_SERIE     linea de serie de la portada
%   PORTADA_ID        identificador de la guia en la portada
%   PORTADA_CIRCULO   el circulito de grado (°), o vacio si no aplica
%   PDF_TITULO        titulo en texto plano (sin \textbf ni comillas TeX) para
%                     los metadatos del PDF (/Title); lo produce lib_xelatex.texto_pdf
%
% Composicion (auditoria 2026-09-03): las bandas de estacion piden espacio
% (\Needspace*) y prohiben el salto justo despues (\nopagebreak), asi no quedan
% huerfanas al pie; las cajas largas (softbox, drawbox, infoband, puentebox) son
% `breakable`, asi una caja de media pagina ya no arrastra todo a la siguiente
% dejando la anterior al 40 %. Todo texto sobre mostaza o verde va en milcNegro
% (blanco sobre #E5B400 da 1,93:1 y sobre #58A923 2,95:1; ninguno pasa WCAG AA).
% El resto de claves las documenta content/guias/_SCHEMA.md. El script valida
% que no queden placeholders sin reemplazar antes de escribir el archivo final.

% Etiquetado PDF (latex-lab): /Lang, estructura y texto alternativo de las
% figuras, para lectores de pantalla. Debe ir ANTES de \documentclass.
\DocumentMetadata{lang=es-CO,tagging=on}
\documentclass[11pt,letterpaper]{article}
% headheight=24pt: la cabecera derecha lleva dos lineas (identidad de la guia
% y estacion actual). headsep se acorta en la misma medida para que el bloque
% de texto quede exactamente donde estaba con headheight=12pt.
\usepackage[margin=2.0cm,headheight=24pt,headsep=13pt]{geometry}
\usepackage[table]{xcolor}
\usepackage{fontspec}
\usepackage[spanish]{babel}
\usepackage{tikz}
% Librerías TikZ para los diagramas de las guías (bloque `recursos:`):
%  · babel  → babel en español vuelve activos caracteres como «>», y TikZ los usa
%             en sus opciones (>=stealth, ->). Sin esta librería, cualquier
%             diagrama con flechas revienta con «\language@active@arg> has an
%             extra }». Debe ir ANTES de que se dibuje nada.
%  · positioning        → sintaxis «below left=2pt and 3pt».
%  · decorations.pathreplacing → llaves ({brace}) para señalar rangos.
\usetikzlibrary{babel,positioning,decorations.pathreplacing}
\usepackage{graphicx}
% Circuitos: el trabajo de electrónica se hace hoy con micro:bit y sus sensores
% integrados (MakeCode, sin cableado), así que no se necesitan esquemáticos.
% Si algún día entran montajes con protoboard, basta descomentar esta línea:
% los diagramas .tex/.tikz declarados en `recursos:` ya se incrustan solos.
% \usepackage{circuitikz}
\usepackage[most]{tcolorbox}
\usepackage{tabularx}
\usepackage{array}
\usepackage{fancyhdr}
\usepackage{titlesec}
\usepackage[document]{ragged2e}  % bandera a la derecha en todo el cuerpo
\usepackage{enumitem}
\usepackage{microtype}
\usepackage{etoolbox}
\usepackage{needspace}
% hyperref al final del preambulo: metadatos del PDF (titulo, autor, idioma,
% familia), marcadores por estacion (\pdfbookmark) y enlaces sin recuadro.
\usepackage[hidelinks,
  pdftitle={El revés del bordado: el error, la frustración y lo que sostiene por detrás},
  pdfauthor={PhD. Álvaro Cárdenas Orozco},
  pdfsubject={Territorio interior · Educación socioemocional · Grado 6° · Momento 6 · MILC v3},
  pdflang={es-CO},
  bookmarksopen=true,bookmarksnumbered=false]{hyperref}

% Helvetica Neue no trae la flecha U+2192, asi que cada «→» escrito literalmente
% en el YAML se compilaba a NADA, en silencio (462 flechas en 61 guias: rutas del
% tipo «paleta → tipografia → lockup» salian rotas). Mapeamos el caracter a su
% version matematica, que si tiene glifo. Asi el autor puede seguir escribiendo
% «→» comodamente en el YAML.
\usepackage{newunicodechar}
\newunicodechar{→}{\ensuremath{\rightarrow}}
% Mismo problema con los simbolos de marca y vinetas: el «✓» de «una palomita ✓»
% salia como cuadrito vacio en el PDF de todas las guias que lo usaban.
\usepackage{pifont}
\newunicodechar{✓}{\ding{51}}
\newunicodechar{✗}{\ding{55}}
\newunicodechar{✔}{\ding{52}}
\newunicodechar{•}{\textbullet}
\newunicodechar{≥}{\ensuremath{\geq}}
\newunicodechar{≤}{\ensuremath{\leq}}

\setmainfont{Helvetica Neue}
\setsansfont{Helvetica Neue}
\setlength{\parindent}{0pt}
\setlength{\parskip}{6pt}
% Sin particion de palabras: en 6.º-7.º una palabra cortada por guion al final
% de linea («Comuni-cacion») frena la lectura mucho mas que una linea corta.
\hyphenpenalty=10000
\exhyphenpenalty=10000
\pretolerance=2000
\tolerance=3000
\setlength{\emergencystretch}{3em}
\linespread{1.10}
\renewcommand{\arraystretch}{1.25}
\setlist{leftmargin=5mm,itemsep=1mm,topsep=1mm}
\newcolumntype{Y}{>{\RaggedRight\arraybackslash}X}

\definecolor{milcMagenta}{HTML}{E6007E}
\definecolor{milcVino}{HTML}{5A0038}
\definecolor{milcTurquesa}{HTML}{0093A5}
\definecolor{milcVerde}{HTML}{58A923}
\definecolor{milcMostaza}{HTML}{E5B400}
\definecolor{milcOcre}{HTML}{B86600}
\definecolor{milcNegro}{HTML}{241816}
\definecolor{milcGris}{HTML}{F7F7F4}
\definecolor{milcLinea}{HTML}{E8E2DD}
\definecolor{milcGradePrimary}{HTML}{0D9488}
\definecolor{milcGradeDark}{HTML}{115E59}
\definecolor{milcGradeSoft}{HTML}{CCFBF1}
\definecolor{milcGradeAccent}{HTML}{CFE7FF}
\definecolor{milcGradeText}{HTML}{FFFFFF}
\color{milcNegro}

\pagestyle{fancy}
\fancyhf{}
% Cabecera de dos lineas: identidad de la guia arriba y, a la derecha y en
% gris, la estacion en curso (\markright desde \milcestacion). Antes de la
% primera estacion la segunda linea va vacia; el \strut mantiene la altura
% para que la linea de arriba no baile entre paginas.
\lhead{\small Territorio interior · Educación socioemocional\\[-1pt]{\footnotesize\strut}}
\rhead{\small Grado 6° · Momento 6 · MILC v3\\[-1pt]{\footnotesize\strut\color{milcNegro!65}\rightmark}}
\cfoot{\small\thepage}
\renewcommand{\headrulewidth}{0pt}

% Sobre mostaza (#E5B400) y verde (#58A923) el texto va en milcNegro; sobre el
% resto de la paleta, en blanco. Se usa dentro de fonttitle/fontupper, donde un
% \color posterior manda sobre coltitle.
\newcommand{\milctintasobre}[1]{%
  \ifboolexpr{test{\ifstrequal{#1}{milcMostaza}} or test{\ifstrequal{#1}{milcVerde}}}%
    {\color{milcNegro}}{\color{white}}}

\newtcolorbox{titlebox}[1]{
  enhanced,colback=#1,colframe=#1,arc=7mm,boxrule=0pt,
  left=7mm,right=7mm,top=3mm,bottom=3mm,
  before skip=8pt,after skip=8pt,
  fontupper=\sffamily\bfseries\Large\color{white}
}

\newtcolorbox{softbox}[3]{
  enhanced,breakable,pad at break=3mm,
  colback=#2,colframe=#1,borderline west={4pt}{0pt}{#1},
  arc=2mm,boxrule=.4pt,left=4mm,right=4mm,top=3mm,bottom=3mm,
  before skip=6pt,after skip=8pt,title={#3},coltitle=white,
  colbacktitle=#1,fonttitle=\sffamily\bfseries\milctintasobre{#1},fontupper=\RaggedRight,
  attach boxed title to top left={xshift=3mm,yshift=-2mm},
  boxed title style={arc=2mm, boxrule=0pt}
}

\newtcolorbox{drawbox}[2]{
  enhanced,breakable,pad at break=4mm,
  colback=white,colframe=#1,arc=4mm,boxrule=1pt,
  left=5mm,right=5mm,top=4mm,bottom=4mm,before skip=8pt,after skip=8pt,
  title={#2},coltitle=white,colbacktitle=#1,fonttitle=\sffamily\bfseries\milctintasobre{#1},
  fontupper=\RaggedRight,
  attach boxed title to top left={xshift=4mm,yshift=-2mm},
  boxed title style={arc=3mm, boxrule=0pt}
}

\newtcolorbox{infoband}[1]{
  enhanced,breakable,pad at break=2mm,colback=#1!8,colframe=#1!50,arc=2mm,boxrule=.5pt,
  left=4mm,right=4mm,top=2mm,bottom=2mm,before skip=4pt,after skip=4pt,
  fontupper=\small\RaggedRight
}

% Puente narrativo entre secciones (contrato editorial Conéctate).
% Da continuidad: "Acabas de X. Ahora vas a Y porque Z."
% Visualmente: banda izquierda en color del grado, texto en itálica pequeña.
\newtcolorbox{puentebox}{
  enhanced,breakable,pad at break=2mm,colback=milcGradeSoft,colframe=milcGradeDark,
  borderline west={3pt}{0pt}{milcGradeDark},
  arc=0mm,boxrule=0pt,left=5mm,right=4mm,top=2.5mm,bottom=2.5mm,
  before skip=4pt,after skip=8pt,
  fontupper=\itshape\small\color{milcNegro}\RaggedRight
}
% La flecha va en modo matematico a proposito: Helvetica Neue NO tiene el glifo
% U+2192, asi que \textrightarrow se compilaba a NADA («Missing character: There
% is no → in font Helvetica Neue») y el puente salia sin flecha. En modo
% matematico la flecha viene de la fuente matematica y si se dibuja.
\newcommand{\puente}[1]{%
  \begin{puentebox}%
    \textcolor{milcGradeDark}{$\rightarrow$}\hspace{2mm}#1%
  \end{puentebox}%
}

\newcommand{\linead}[1]{\noindent\color{milcLinea}\rule{#1}{0.5pt}\color{milcNegro}}
\newcommand{\respuesta}[1]{\par\vspace{2mm}\foreach \n in {1,...,#1}{\noindent\color{milcLinea}\rule{\linewidth}{0.5pt}\par\vspace{5mm}}\color{milcNegro}}
\newcommand{\checkbox}{\textcolor{milcGradeDark}{\fbox{\phantom{x}}}\hspace{5pt}}

% ─── Listas aireadas para las guías socioemocionales ────────────────────
% Componer las verificaciones como «\checkbox texto\\» deja el texto que salta
% de línea debajo del cuadrito y apelmaza el bloque. Estos entornos alinean la
% continuación y airean los ítems. Los usa Territorio Interior; cualquier
% familia puede hacerlo.
\newlist{checklista}{itemize}{1}
\setlist[checklista]{
  label=\textcolor{milcGradeDark}{\fbox{\phantom{x}}},
  leftmargin=9mm, labelsep=3.2mm, itemsep=2.4mm,
  topsep=2.5mm, parsep=0pt, partopsep=0pt, before=\RaggedRight
}
\newlist{pasoslista}{enumerate}{1}
\setlist[pasoslista]{
  label=\textcolor{milcGradeDark}{\sffamily\bfseries\arabic*.},
  leftmargin=8mm, labelsep=3mm, itemsep=2.8mm,
  topsep=1mm, parsep=0pt, partopsep=0pt, before=\RaggedRight
}

\newcommand{\phasecard}[4]{
\begin{tcolorbox}[enhanced,colback=#3,colframe=#2,arc=3mm,boxrule=.5pt,height=3.05cm,valign=center,halign=center,left=1.5mm,right=1.5mm,top=2mm,bottom=2mm]
{\sffamily\bfseries\fontsize{22}{24}\selectfont\textcolor{#2}{#1}}\par
{\sffamily\bfseries\small #4}
\end{tcolorbox}
}

% ── Piezas del rediseno 2026 (legibilidad 6.º-11.º) ───────────────────
% Banda de estacion: reemplaza el titlebox macizo. Titulo a la izquierda,
% dato practico (tiempo/modalidad) a la derecha, en una sola linea.
% Nunca huerfana: pide 9 lineas libres (o salta de pagina rellenando con
% \vfill) y prohibe el salto justo despues de la banda. Nueve y no seis:
% la caja partible que sigue necesita ~80pt (titulo adosado, relleno y dos
% lineas) para arrancar en la misma pagina; con menos, tcolorbox la manda
% entera a la siguiente y la banda se queda sola. El penalty 10000 va
% en `after`, ANTES del espacio posterior: un `after skip` (aunque sea 0pt)
% es un \vskip tras la caja, y ese glue es un punto de corte legal que un
% \nopagebreak escrito despues de \end{tcolorbox} ya no cierra. Cada banda
% deja marcador PDF y actualiza la cabecera derecha.
\newcounter{milcestacion}
\newcommand{\milcestacion}[2]{%
\Needspace*{9\baselineskip}%
\stepcounter{milcestacion}%
\pdfbookmark[1]{#2}{milcestacion\themilcestacion}%
\markright{#2}%
\begin{tcolorbox}[enhanced,colback=#1,colframe=#1,arc=2mm,boxrule=0pt,
  left=5mm,right=5mm,top=2.6mm,bottom=2.6mm,before skip=11pt,
  after={\par\nopagebreak\vskip7pt}]
{\sffamily\bfseries\fontsize{15}{18}\selectfont\milctintasobre{#1}#2}
\end{tcolorbox}}

% Tarjeta de paso numerada: sustituye las tablas «Pilar / Aplicacion», que
% eran formato de consulta docente, no de estudiante. El numero va en un
% circulo sobre el borde para que la secuencia se lea de un vistazo.
\newcommand{\milcpaso}[3]{%
\Needspace{4\baselineskip}%
\begin{tcolorbox}[enhanced,colback=white,colframe=milcLinea,arc=1.5mm,boxrule=.9pt,
  left=14mm,right=4mm,top=3mm,bottom=3mm,before skip=4pt,after skip=4pt,
  fontupper=\RaggedRight,
  overlay={\node[anchor=north west,circle,fill=#2,text=white,inner sep=0pt,
    minimum size=7.4mm,font=\sffamily\bfseries\small]
    at ([xshift=3.6mm,yshift=-3.2mm]frame.north west) {#1};}]
#3
\end{tcolorbox}}

% Casilla marcable de tamano util con lapiz (la anterior era \fbox{\phantom{x}},
% de alto variable segun el texto que la seguia).
\newcommand{\milcmarca}{\framebox[3.8mm]{\rule{0pt}{3.4mm}}}

% Fila marcable de lista suelta (checks de la estacion 1).
\newcommand{\milccheckrow}[1]{%
\par\vspace{1.8mm}\noindent
\makebox[6.4mm][l]{\raisebox{-.55mm}{\textcolor{milcNegro}{\milcmarca}}}%
\parbox[t]{\dimexpr\linewidth-6.4mm\relax}{\RaggedRight #1}\par}

% Celda marcable dentro de la rubrica (tabla de autoevaluacion). Misma
% sangria francesa, pero pensada para vivir dentro de una columna X.
\newcommand{\milcopcion}[1]{%
\makebox[5.2mm][l]{\raisebox{-.55mm}{\textcolor{milcNegro}{\milcmarca}}}%
\parbox[t]{\dimexpr\linewidth-5.2mm\relax}{\RaggedRight #1}}

% ── Recursos visuales de la guía ──────────────────────────────────────
% Declarados en el bloque `recursos:` del YAML y emitidos por el builder.
% \guiaFigura   → imagen rasterizada o PDF (foto, carta celeste, montaje).
% \guiaDiagrama → fuente TikZ/circuitikz (.tex) incrustada como vector:
%                 es la vía para circuitos y esquemáticos editables.
% [#1] = texto alternativo (alt) · #2 = ruta relativa al .tex · #3 = pie
% El alt llega al PDF etiquetado (/Alt) y lo lee el lector de pantalla.
\newcommand{\guiaFigura}[3][]{%
\begin{center}
\includegraphics[width=0.88\linewidth,height=0.38\textheight,keepaspectratio,alt={#1}]{#2}\\[1.5mm]
{\footnotesize\itshape\textcolor{milcNegro!75}{#3}}
\end{center}
}

\newcommand{\guiaDiagrama}[2]{%
\begin{center}
\input{#1}\\[1.5mm]
{\footnotesize\itshape\textcolor{milcNegro!75}{#2}}
\end{center}
}

\begin{document}
\thispagestyle{empty}

% ── Portada ───────────────────────────────────────────────────────────
% Antes: un rectangulo de color a pagina completa. Se veia bien en pantalla,
% pero en la fotocopiadora del colegio se comia el toner y salia gris sucio.
% Ahora la banda ocupa el tercio superior y el resto va en blanco.
\begin{tikzpicture}[remember picture,overlay]
  \path[fill=milcGradePrimary]
    (current page.north west) rectangle ([yshift=-10.6cm]current page.north east);

  \node[anchor=north west,text=milcGradeText] at ([xshift=2.0cm,yshift=-1.55cm]current page.north west)
    {\sffamily\bfseries\fontsize{12}{14}\selectfont MOMENTO};
  \node[anchor=north west,text=milcGradeText,opacity=.85] at ([xshift=2.0cm,yshift=-2.35cm]current page.north west)
    {\sffamily\bfseries\fontsize{8.5}{10}\selectfont TERRITORIO INTERIOR · SOCIOEMOCIONAL};
  \node[anchor=north east,text=milcGradeText,opacity=.85,align=right] at ([xshift=-2.0cm,yshift=-1.55cm]current page.north east)
    {\sffamily\bfseries\fontsize{9}{12}\selectfont I.E. Sor María Juliana\\Cartago, Valle del Cauca};

  \node[anchor=west,text=milcGradeText] (gradonum) at ([xshift=1.85cm,yshift=-7.1cm]current page.north west)
    {\sffamily\bfseries\fontsize{112}{112}\selectfont 6};
  % El circulito de grado (°) solo aplica a las guias de grado («11°»); en
  % Bebras y Semillero el numero grande es un momento/modulo, no un grado.
  % Se posiciona relativo al nodo `gradonum`, no en coordenadas absolutas.
  

  \node[anchor=west,text=milcGradeText,text width=11.4cm,align=left] at ([xshift=1.5cm,yshift=0.5cm]gradonum.east)
    {
     {\sffamily\bfseries\fontsize{9}{11}\selectfont\MakeUppercase{Territorio interior · Socioemocional}}\\[3mm]
     {\sffamily\bfseries\fontsize{21}{25}\selectfont\setlength{\baselineskip}{27pt}El revés del bordado\\[4pt]el error, la frustración\\[4pt]y lo que sostiene}\\[3.5mm]
     {\sffamily\bfseries\fontsize{15}{18}\selectfont Grado 6° · Momento 6}
    };
\end{tikzpicture}

\vspace*{9.4cm}
\vfill

{\sffamily\bfseries\fontsize{8}{10}\selectfont\textcolor{milcGradeDark}{LO QUE TE LLEVAS HOY}}

\begin{tcolorbox}[enhanced,colback=white,colframe=milcNegro,arc=2mm,boxrule=1.6pt,
  left=5mm,right=5mm,top=4mm,bottom=4mm,before skip=3pt,after skip=12pt,
  fontupper=\RaggedRight\fontsize{11}{15.5}\selectfont]
Tu revés: la hoja donde muestras el proceso de algo que sí lograste —los intentos tachados, lo que cada uno enseñó y el remiendo que hiciste—, más una cosa concreta que vas a rehacer esta semana.
\end{tcolorbox}

\begin{tcolorbox}[enhanced,colback=milcGris,colframe=milcGris,arc=2mm,boxrule=0pt,
  left=5mm,right=5mm,top=3.5mm,bottom=3.5mm,after skip=12pt]
{\sffamily\bfseries\fontsize{8}{10}\selectfont\textcolor{milcNegro!55}{ESTA GUÍA ES DE}}\\[3mm]
\begin{tabularx}{\linewidth}{X p{3.2cm} p{3.2cm}}
\linead{\linewidth} & \linead{3.0cm} & \linead{3.0cm}\\[-1mm]
{\fontsize{8.5}{10}\selectfont\textcolor{milcNegro!55}{Nombre completo}} &
{\fontsize{8.5}{10}\selectfont\textcolor{milcNegro!55}{Grupo}} &
{\fontsize{8.5}{10}\selectfont\textcolor{milcNegro!55}{Fecha}}\\
\end{tabularx}
\end{tcolorbox}

\vfill

\noindent\textcolor{milcLinea}{\rule{\linewidth}{1pt}}\\[2mm]
{\sffamily\bfseries\fontsize{9.5}{12}\selectfont Autor: Dr. Álvaro Cárdenas Orozco}\\[1mm]
{\sffamily\fontsize{8.5}{11}\selectfont\textcolor{milcNegro!55}{Metodología MILC · Apertura ancestral · Error y frustración · Triángulo Dussel-Estoicismo-Floridi}}

\newpage

\section*{Apertura: El revés del bordado: el error, la frustración y lo que sostiene por detrás}

\begin{softbox}{milcGradeDark}{milcGradeSoft}{Saber ancestral}
Voltea un bordado. \textbf{Por el derecho está el diseño; por el revés están los nudos}, las puntas sueltas, los hilos que cruzan de un lado a otro. Nadie enmarca el revés, y sin embargo \textbf{es el revés el que sostiene el derecho}: si sueltas los nudos, el dibujo se deshace. En el calado cartagüeño esto se vuelve todavía más claro. Cuando una caladora saca un hilo que no debía ---y pasa---, la simetría de la rejilla se daña, y entonces \textbf{remienda}: mueve los hilos para cerrar el espacio de más, o teje uno nuevo en su lugar. El remiendo tiene una regla, y la dijo Olivia, caladora de 65 años: \emph{«no quieres que se vea remendada»}. Es decir, \textbf{el arreglo está hecho para ser invisible} (Pérez-Bustos, 2019). Y hay un límite honesto: cuando el daño es demasiado grande ---demasiados hilos retirados donde debía haber menos---, \textbf{ya no hay remiendo posible}, y también hay que saber reconocerlo. \emph{Tres cosas que te sirven hoy: los nudos van por detrás, el arreglo casi siempre existe, y a veces toca empezar de nuevo.}
\end{softbox}

\begin{softbox}{milcTurquesa}{milcGris}{Saber contemporáneo}
En el colegio pasa algo raro: solo se muestra el derecho. Se entrega la hoja limpia, la respuesta correcta, el trabajo pasado en limpio ---y los borradores, los tachones y los tres intentos anteriores van a la basura---. Así, cada quien ve \textbf{su propio revés y el derecho de los demás}, y llega a una conclusión falsa: \emph{«a todos les sale a la primera menos a mí»}. Es mentira, y es una mentira que hace daño. La investigación educativa lo ha medido: Manu Kapur mostró en 2008 que los estudiantes que \textbf{primero pelean con un problema difícil} ---y fallan--- y \emph{después} reciben la explicación, aprenden más hondo que los que reciben la explicación de entrada. A eso lo llamó \textbf{fracaso productivo} (Kapur, 2008). \emph{No es que el error sea bonito, ni que haya que celebrarlo:} es que la pelea con lo que no sale prepara la cabeza para entender. \textbf{Y hay una distinción que a los once años cambia el año entero:} \emph{equivocarse} es algo que uno hace; \emph{ser un fracaso} es algo que uno cree que es. La primera se arregla. La segunda ni siquiera es cierta.
\end{softbox}

\begin{softbox}{milcMostaza}{milcGris}{Pregunta-puente}
\textit{La caladora remienda para que el arreglo no se vea, pero \textbf{sabe perfectamente dónde está} ---ella se acuerda de cada nudo---. \textbf{¿Qué has aprendido a hacer bien que hoy parece fácil}\ldots y cuántos intentos tachados hubo por detrás que nadie vio?}
\end{softbox}

\begin{softbox}{milcVerde}{milcGris}{Saber y saber hacer}
Al terminar podrás: (1) \textbf{identificar} la diferencia entre equivocarte y creerte un fracaso, en casos tuyos de verdad; (2) \textbf{explicar} por qué pelear con algo difícil antes de que te lo expliquen ayuda a aprender, y qué hacer con la frustración cuando aparece; (3) \textbf{crear} la hoja de tu revés: el proceso completo de algo que sí lograste; (4) \textbf{decidir} un remiendo concreto ---algo que vas a rehacer esta semana--- y distinguirlo de lo que ya no tiene arreglo y hay que empezar otra vez.
\end{softbox}



\small
\begin{tabularx}{\textwidth}{>{\bfseries}p{3.0cm}Y}
\rowcolor{milcGradePrimary!12}Autor & Dr. Álvaro Cárdenas Orozco\\
DBA & Comprende los fenómenos que afectan a su territorio, regula sus emociones ante ellos y acompaña a otros con empatía (transversal 6°--11°, articulado con la política «Escuela, territorio de vida» del MEN, Resolución 006519 de 2025, y la Ley 1620 de 2013)\\
Referentes & Primera ayuda psicológica (OMS, 2011/2012) · Directrices IASC (2007) · USGS y Servicio Geológico Colombiano · Pueblo nasa y la minga (CRIC) · Dussel · Estoicismo · Floridi\\
Producto final & Tu revés: la hoja donde muestras el proceso de algo que sí lograste —los intentos tachados, lo que cada uno enseñó y el remiendo que hiciste—, más una cosa concreta que vas a rehacer esta semana.\\
\end{tabularx}
\normalsize

\puente{Ya tienes palabras para lo que sientes, un fogón para no reaccionar de una y una técnica para respirar. Hoy usas las tres en el momento más común de todos: \textbf{cuando algo no te sale}.}

\milcestacion{milcGradeDark}{Ruta MILC: así vamos a aprender}
\begin{tabularx}{\textwidth}{>{\centering\arraybackslash}X>{\centering\arraybackslash}X>{\centering\arraybackslash}X>{\centering\arraybackslash}X}
\phasecard{1}{milcVerde}{milcGris}{Escucha\\[-1mm]\small Observo, escucho y pregunto.} &
\phasecard{2}{milcTurquesa}{milcGris}{Sistematización\\[-1mm]\small Distingo y sostengo.} &
\phasecard{3}{milcMagenta}{milcGris}{Praxis\\[-1mm]\small Construyo y aplico.} &
\phasecard{4}{milcVino}{milcGris}{Evaluación\\[-1mm]\small Reflexión liberadora.}
\end{tabularx}


\puente{Primero mira tus propios reveses, sin adornarlos. No es una lista de defectos: es la parte del trabajo que nadie muestra.}

\milcestacion{milcVerde}{Estación 1 de 3 · Escucha}

\begin{softbox}{milcVerde}{milcGris}{Escena para escuchar}
\textbf{Actividad 1 · IDENTIFICA --- Dos listas, en privado.} \textbf{Nadie va a leer esta página.} \textbf{Lista A (5 min).} Escribe \textbf{tres cosas que hoy haces bien} ---leer en voz alta, montar en bicicleta, cocinar algo, un juego, un deporte, dibujar, cuidar a alguien menor---. Al lado de cada una, contesta de memoria: \emph{¿te salió bien la primera vez?} y \emph{¿cuántas veces la hiciste mal antes?}. \textbf{Lista B (5 min).} Escribe \textbf{dos cosas que dejaste de intentar} porque no te salían. Al lado, la frase exacta que te dijiste cuando decidiste dejarlas ---escríbela tal cual, aunque sea fea: «yo soy malo para esto», «no nací para eso», «me da pena»---. \textbf{Y una línea final:} mira la Lista A y la Lista B juntas. \emph{¿Qué tenían de distinto las tres primeras?} Casi siempre la respuesta no es el talento: es que en esas nadie te vio equivocarte, o no te importó que te vieran.
\end{softbox}

{\sffamily\bfseries\fontsize{8}{10}\selectfont\textcolor{milcNegro!55}{MARCA CUANDO LO HAYAS HECHO}}
\milccheckrow{Escribiste tres cosas que haces bien, con las veces que te salieron mal antes.}
\milccheckrow{Escribiste dos cosas que dejaste de intentar, con la frase exacta que te dijiste.}
\milccheckrow{Comparaste las dos listas y anotaste qué tenían de distinto las de la Lista A.}

\begin{infoband}{milcVerde}


\textbf{En tu cuaderno (Actividad 1 · IDENTIFICA).} Título: «Actividad 1. Mi derecho y mi revés». \textbf{Formato:} dos listas ---tres cosas que hago bien con sus intentos fallidos, dos que abandoné con la frase que me dije--- y una línea de comparación. \textbf{Extensión:} media página. \textbf{Sabes que terminaste cuando} escribiste esa frase fea \textbf{tal como te la dijiste}, sin suavizarla. Esta página es tuya: \textbf{nadie más tiene que leerla si tú no quieres}.
\end{infoband}

\puente{Ya los tienes a la vista. Ahora entiende qué es equivocarse, qué es la frustración y por qué pelear con lo difícil no es perder el tiempo.}

\milcestacion{milcTurquesa}{Fase 2 · Sistematización: entender el error}

\begin{softbox}{milcTurquesa}{milcGris}{Cuatro claves sobre equivocarse}
Cuatro claves para que equivocarse deje de ser una emergencia y vuelva a ser lo que siempre fue: la parte de atrás del trabajo.
\end{softbox}

\milcpaso{1}{milcTurquesa}{\textbf{Equivocarse es algo que haces; ser un fracaso es algo que crees que eres.} La diferencia parece de palabras y es enorme. \emph{«Me equivoqué en el ejercicio»} deja algo por hacer: revisar dónde. \emph{«Soy malo para matemáticas»} no deja nada, porque no describe un error sino una identidad ---y contra una identidad no se puede trabajar---. \textbf{Truco para pillarte:} si la frase empieza con «soy», casi siempre es la segunda; si empieza con «me salió» o «hice», es la primera. \emph{Cambiar una por otra no es autoengaño: es volver a poner el problema donde sí se puede resolver.}}
\milcpaso{2}{milcTurquesa}{\textbf{La frustración es una emoción, y estaba en tu muestrario.} Aparece cuando quieres algo y algo se interpone; es prima de la rabia y trae su propio impulso: \textbf{tirar el cuaderno, cerrar el computador, decir «no me importa»}. Y aquí ya sabes qué hacer, porque es exactamente lo del momento 3: \textbf{la emoción informa, no obliga}. Nómbrala (\emph{«esto es frustración y va en 7»}), haz tu pausa, respira con la salida larga como en el momento 4, y \textbf{vuelve al intento} ---o decide parar hoy y volver mañana, que también es una decisión y no una rendición---.}
\milcpaso{3}{milcTurquesa}{\textbf{Pelear con lo difícil antes de que te lo expliquen te hace aprender más.} No es una frase de cartel: se midió. Kapur (2008) puso a unos estudiantes a pelear primero con problemas difíciles, sin ayuda, y a otros a recibir la explicación de una. \textbf{Los que primero fallaron aprendieron más hondo}, y él lo llamó \emph{fracaso productivo}. La razón es sencilla: mientras peleas, tu cabeza descubre \textbf{dónde está la dificultad}, y cuando llega la explicación ya sabe dónde ponerla. \emph{Por eso rendirse en el minuto dos no es ahorrar tiempo: es saltarse la parte que te iba a servir.}}
\milcpaso{4}{milcTurquesa}{\textbf{Hay daños con remiendo y daños sin remiendo, y saber cuál es cuál es parte del oficio.} La caladora remienda casi todo, pero cuando faltan demasiados hilos \textbf{no insiste}: reconoce que ahí ya no hay arreglo (Pérez-Bustos, 2019). En tu vida escolar pasa igual: una tarea mal hecha se rehace, un examen malo se compensa, un trabajo se corrige. Pero hay cosas que no se remiendan y \textbf{se empiezan de nuevo} ---el proyecto planteado desde una idea equivocada, la amistad que ya no era---. \emph{Distinguirlas ahorra semanas de estar cosiendo algo que ya no se sostiene.}}

\begin{drawbox}{milcTurquesa}{Qué hacer cuando algo no sale, en orden}
\begin{pasoslista}
\item \textbf{Nombra lo que sientes.} Frustración, rabia, vergüenza o las tres. Con la palabra exacta, como en el momento 2.
\item \textbf{Haz tu pausa.} La del momento 3, o el aire del momento 4. Treinta segundos alcanzan para no tirar el cuaderno.
\item \textbf{Cambia la frase.} De «soy malo para esto» a «esto todavía no me sale». Fíjate en la palabra \emph{todavía}: es la que deja la puerta abierta.
\item \textbf{Busca dónde exactamente falló.} No «todo mal»: \emph{cuál paso}. Casi siempre es uno solo, y casi nunca es el que uno creía.
\item \textbf{Decide: remiendo o empezar de nuevo.} Las dos son válidas. Lo que no sirve es quedarse cosiendo lo que ya no se sostiene.
\end{pasoslista}
\end{drawbox}

\begin{softbox}{milcMostaza}{milcGris}{Errores comunes y su corrección}
\textbf{Tirar el trabajo entero por un paso malo:} casi siempre falló uno, no todo. \textbf{Esconder el proceso:} si borras los intentos, te quedas sin la evidencia de que avanzaste. \textbf{Confundir «no me sale» con «no puedo»:} entre las dos está la palabra \emph{todavía}. \textbf{Insistir donde ya no hay remiendo:} volver a empezar no es fracasar, es un juicio técnico. \textbf{Compararte con el derecho de los demás:} tú ves tu revés completo y de ellos solo la cara bonita. \textbf{Pedir ayuda solo cuando ya te rendiste:} pedirla en la mitad del intento es más útil y no es hacer trampa.
\end{softbox}

\begin{infoband}{milcTurquesa}


\textbf{En tu cuaderno (Actividad 2 · EXPLICA).} Título: «Actividad 2. Qué hacer cuando no sale». \textbf{Formato:} los cinco pasos del recuadro aplicados a \textbf{una} de las dos cosas que abandonaste (Lista B), y debajo la frase vieja reescrita con «todavía». \textbf{Extensión:} media página. \textbf{Sabes que terminaste cuando} puedes señalar \textbf{el paso exacto} donde se te atoraba, en vez de decir «no me entraba».
\end{infoband}

\puente{Ya lo entiendes. Es hora de hacer algo que casi nunca se hace en un colegio: \textbf{mostrar el revés}, y mostrarlo con orgullo.}

\milcestacion{milcMagenta}{Fase 3 · Praxis: mostrar el revés}

\begin{softbox}{milcMagenta}{milcGris}{Cómo se arma la hoja del revés}
Hoy se hace lo contrario de lo que hace un colegio: \textbf{se muestra el revés}. Vas a armar una hoja donde se vea el proceso completo de algo que sí lograste ---con tachones y todo--- porque esa hoja es la prueba de que sabes trabajar, no solo de que llegaste.
\end{softbox}

\milcpaso{1}{milcMagenta}{\textbf{Elegir el logro (5 min).} De tu Lista A, escoge \textbf{uno} del que te acuerdes bien del proceso: aprender a montar en bicicleta, un dibujo que repetiste, una jugada, una receta, un trabajo del año pasado. Que no sea el más impresionante: que sea del que \textbf{te acuerdes de los intentos}.}
\milcpaso{2}{milcMagenta}{\textbf{Recuperar los intentos (12 min).} Escribe o dibuja, en orden, \textbf{cómo fue cada intento} y qué salió mal en cada uno. Tres como mínimo. Si conservas los borradores o las fotos, pégalos. Si no, reconstrúyelos de memoria ---sirve igual---. \emph{Ojo con la trampa de la memoria: casi siempre recordamos el resultado y borramos el proceso, así que este paso cuesta más de lo que parece.}}
\milcpaso{3}{milcMagenta}{\textbf{Nombrar lo que enseñó cada intento (10 min).} Al lado de cada uno, una línea: \emph{¿qué supe después de este intento que no sabía antes?}. No vale «que estaba mal»: vale «que había que frenar antes de la curva», «que la mezcla iba más espesa», «que yo estaba leyendo el enunciado al revés». \textbf{Esa columna es el verdadero producto de hoy.}}
\milcpaso{4}{milcMagenta}{\textbf{Escoger el remiendo (8 min).} Ahora mira tu Lista B ---lo que abandonaste--- y elige \textbf{una} cosa que vas a rehacer esta semana. Concreta y pequeña: no «volver a ser bueno en matemáticas», sino «rehacer los tres ejercicios que me salieron mal, uno por día, y preguntar el jueves». \emph{Y si al mirarla decides que esa no tiene remiendo y prefieres empezar otra cosa, escríbelo también: esa decisión vale igual.}}

\begin{drawbox}{milcMagenta}{Antes de cerrar tu revés}
\begin{checklista}
\item Tu hoja muestra \textbf{al menos tres intentos} y qué salió mal en cada uno.
\item Cada intento tiene su línea de «qué me enseñó», y ninguna dice solo «que estaba mal».
\item Reescribiste la frase fea de la Actividad 1 usando la palabra \textbf{todavía}.
\item Tu remiendo es concreto, pequeño y tiene día: se puede hacer esta semana.
\item Sabes decir cuál de las dos cosas que abandonaste \textbf{sí} tiene arreglo y cuál no, con una razón.
\item Estás dispuesto a que alguien vea esta hoja, o decidiste a conciencia que es solo tuya. Las dos valen.
\end{checklista}
\end{drawbox}

\begin{softbox}{milcMagenta}{milcGris}{Plantilla de trabajo}
\textbf{Plantilla de la hoja del revés.} \emph{Lo que hoy hago bien:} \_\_\_\_ · \emph{Intento 1:} qué hice / qué salió mal / qué aprendí · \emph{Intento 2:} \ldots · \emph{Intento 3:} \ldots · \emph{Lo que finalmente funcionó:} \_\_\_\_ \\[1mm]
\textbf{Y el remiendo.} \emph{«Voy a rehacer \_\_\_\_, el día \_\_\_\_, y sabré que quedó cuando \_\_\_\_».} \\[1mm]
\textbf{La frase, corregida.} \emph{Antes:} «\_\_\_\_». \emph{Ahora:} «\_\_\_\_ todavía no me sale».
\end{softbox}

\begin{infoband}{milcMagenta}


\textbf{En tu cuaderno (Actividad 3 · CREA).} Título: «Actividad 3. La hoja de mi revés». \textbf{Formato:} el logro elegido, los tres intentos con su columna de «qué me enseñó», la frase reescrita y el remiendo con día. \textbf{Extensión:} una página. \textbf{Sabes que terminaste cuando} tu hoja le serviría a alguien que \textbf{está empezando} eso mismo ---porque ahí están los errores que va a cometer---.
\end{infoband}

\puente{El producto es esa hoja. \emph{Un trabajo sin proceso visible no demuestra que aprendiste: demuestra que llegaste.}}

\milcestacion{milcMostaza}{Producto de la sesión}
\textbf{La hoja del revés: el proceso completo de un logro, más un remiendo con fecha}.
\begin{softbox}{milcMostaza}{milcGris}{Criterios de calidad}
(1) Se muestran al menos tres intentos fallidos, en orden, con lo que salió mal en cada uno. (2) Cada intento tiene escrito qué enseñó, en términos concretos y no en «estaba mal». (3) La frase de identidad («soy malo para…») está reescrita como algo que todavía no sale. (4) El remiendo es concreto, pequeño, tiene día asignado y una forma de saber que quedó hecho. (5) Se distingue, con una razón, lo que tiene arreglo de lo que conviene empezar de nuevo.
\end{softbox}

\puente{Antes de cerrar, separa bien las dos cosas: lo que tiene remiendo y lo que ya no. Confundirlas cuesta semanas.}

\milcestacion{milcVino}{Evaluación liberadora · cinco dimensiones}

\begin{softbox}{milcVino}{milcGris}{Sentido de la evaluación}
Evaluar no es solo revisar una nota. Es reconocer qué aprendí, qué cambió en mí, qué emociones manejé, cómo me relaciono con la ciudad y la comunidad, y qué le dejaré a quien viene detrás.
\end{softbox}

\textbf{Mi reflexión liberadora en cinco dimensiones.} Completa cada frase iniciada:

\small
\begin{tabularx}{\textwidth}{>{\bfseries\RaggedRight\arraybackslash}p{3.4cm}Y>{\RaggedRight\arraybackslash}p{4.6cm}}
\rowcolor{milcVino}\color{white}Dimensión & \color{white}\bfseries Mi respuesta (frame starter) & \color{white}\bfseries Algo que descubrí\\
1. Personal & Lo que más cambió en mí fue \linead{6.5cm} & \linead{4.2cm}\\
2. Emocional & La emoción más fuerte fue \linead{2.5cm} y la regulé \linead{3cm} & \linead{4.2cm}\\
3. Ciudadana & En democracia esto importa porque \linead{6cm} & \linead{4.2cm}\\
4. Local (Cartago) & En mi barrio sería útil para \linead{6.5cm} & \linead{4.2cm}\\
5. Intergeneracional & A mi abuela/abuelo le diría \linead{6.5cm} & \linead{4.2cm}\\
\end{tabularx}
\normalsize

\begin{infoband}{milcVino}
\textbf{Para la próxima sustentación.} Vuelve a esta tabla en seis meses. Verás que las respuestas cambian — no porque lo aprendido cambie, sino porque tú cambias. La evaluación liberadora es una conversación contigo a través del tiempo.
\end{infoband}


\puente{Cerramos con tres voces: el trabajo invisible que sostiene lo visible, cambiar de opinión cuando te demuestran que estabas equivocado, y qué significa cultivar en vez de solo conservar.}

\milcestacion{milcGradeDark}{Triángulo de pensamiento · cierre filosófico}

\begin{softbox}{milcVino}{milcGris}{Dussel · lente del nosotros}
\textit{«El otro se revela realmente como otro\ldots como el pobre, el oprimido; el que, a la vera del camino, fuera del sistema, muestra su rostro sufriente y sin embargo desafiante: «¡Tengo hambre!, ¡tengo derecho a comer!»»}

{\footnotesize\itshape\color{milcNegro!70}--- Enrique Dussel · Filosofía de la liberación (1977)}

El remiendo de la caladora está hecho para no verse, y eso tiene un costo: \textbf{el trabajo que no se ve, no se paga ni se reconoce}. Las caladoras lo dicen ---«la gente no paga por este tipo de trabajo»--- y les pasa lo que Dussel describe: quedan \emph{fuera del sistema} justamente por hacer bien la parte invisible. \textbf{Vale para tu curso:} el que se demora tres horas en lo que otro hace en veinte minutos no es peor; está haciendo un revés que nadie mira. \emph{Preguntar «¿cuánto le costó?» antes de decir «qué fácil le queda» es un acto de justicia pequeño y diario.}

\textbf{¿A quién de mi curso he juzgado por su derecho sin saber nada de su revés?}
\end{softbox}
\linead{\linewidth}\\[1mm]
\linead{\linewidth}

\begin{softbox}{milcMostaza}{milcGris}{Estoicismo · Marco Aurelio · Meditaciones VI, 21 (c. 175 d.C.) · cuidado interior}
\textit{«Si alguien puede refutarme y probar de modo concluyente que pienso o actúo incorrectamente, de buen grado cambiaré de proceder. Pues persigo la verdad, que no dañó nunca a nadie; en cambio, sí se daña el que persiste en su propio engaño e ignorancia.»}

{\footnotesize\itshape\color{milcNegro!70}--- Marco Aurelio · Meditaciones VI, 21 (c. 175 d.C.)}

Un emperador escribiendo para sí mismo que \textbf{cambiará de proceder de buen grado} si alguien le demuestra que está equivocado. Fíjate en lo que dice al final: el daño no está en equivocarse, \textbf{está en persistir}. Cuando alguien te corrige un ejercicio o te dice que tu parte del trabajo quedó mal, ahí hay dos caminos: defenderte ---y quedarte igual--- o cambiar de proceder ---y quedar mejor---. \emph{La corrección se siente como un ataque y casi nunca lo es: es información que llegó gratis.}

\textbf{La última vez que me corrigieron, ¿me defendí o revisé? ¿Qué habría pasado si hacía lo otro?}
\end{softbox}
\linead{\linewidth}\\[1mm]
\linead{\linewidth}

\begin{softbox}{milcTurquesa}{milcGris}{Floridi · lente de la infoesfera}
\textit{«El florecimiento de las entidades informacionales y de toda la infoesfera debe promoverse preservando, cultivando y enriqueciendo sus propiedades.»}

{\footnotesize\itshape\color{milcNegro!70}--- Luciano Floridi · La ética de la información (2013; trad.)}

Floridi no dice solo \emph{preservar}: dice \textbf{cultivar y enriquecer}, que es dejar las cosas mejor de como estaban. Tu hoja del revés hace precisamente eso: convierte tus errores ---que se iban a la basura--- en algo que le sirve al que viene detrás. \textbf{Un error tirado se pierde; un error escrito se vuelve mapa.} Por eso los oficios que duran ---el calado, la carpintería, la medicina--- no guardan solo los resultados: guardan los arreglos, las fallas y lo que se aprendió de ellas.

\textbf{De los errores que cometí este mes, ¿cuál le serviría a alguien que va a empezar lo mismo?}
\end{softbox}
\linead{\linewidth}\\[1mm]
\linead{\linewidth}

\begin{infoband}{milcNegro}
\textbf{Nota para el docente.} Las tres son citas textuales y llevan su referencia debajo. Puedes remitir a la obra en clase.
\end{infoband}

\milcestacion{milcVino}{Mi autoevaluación · marca una casilla por fila}

{\small
\setlength{\extrarowheight}{2.2pt}
\begin{tabularx}{\textwidth}{>{\RaggedRight\arraybackslash\bfseries}p{3.0cm}YYY}
\rowcolor{milcVino}\color{white}\bfseries Criterio & \color{white}\bfseries Lo logré & \color{white}\bfseries En proceso & \color{white}\bfseries Necesito apoyo\\[.6mm]
Comprensión del tema & \milcopcion{Explico las ideas con mis palabras.} & \milcopcion{Reconozco las ideas, falta precisión.} & \milcopcion{Necesito repasar lo básico.}\\[.8mm]
\rowcolor{milcGris}Manejo del error & \milcopcion{Distingo equivocarme de ser un fracaso, y sé qué me enseñó cada intento.} & \milcopcion{Reconozco el error, pero todavía me quedo en «soy malo para esto».} & \milcopcion{Cuando algo no me sale, dejo de intentarlo.}\\[.8mm]
Aplicación y producto & \milcopcion{Mi producto cumple los criterios.} & \milcopcion{Está armado, con ajustes pendientes.} & \milcopcion{Aún está en construcción.}\\[.8mm]
\rowcolor{milcGris}Reflexión 5 dimensiones & \milcopcion{Respondí cada una con honestidad.} & \milcopcion{Respondí algunas con detalle.} & \milcopcion{Mis respuestas son muy generales.}\\[.8mm]
Triángulo de pensamiento & \milcopcion{Conecté las 3 voces con mi trabajo.} & \milcopcion{Conecté al menos una.} & \milcopcion{No las relaciono con mi proceso.}\\[.8mm]
\end{tabularx}}

\vspace{3mm}
\textbf{Hoy entendí que:}\\[1mm]
\linead{\linewidth}\\[1mm]
\linead{\linewidth}

\textbf{Mi compromiso para la próxima sesión es:}\\[1mm]
\linead{\linewidth}\\[1mm]
\linead{\linewidth}

\begin{softbox}{milcMostaza}{milcGris}{Cita de despedida · Marco Aurelio}
\textit{``El obstáculo en el camino se vuelve el camino. Lo que estorba al camino se vuelve camino.''}\\[1mm]
Cada error de hoy, bien observado, se convierte en pieza del camino que pisarás mañana.
\end{softbox}

\small
\begin{tabularx}{\textwidth}{>{\bfseries}p{3.6cm}Y}
\rowcolor{milcGradeDark!12}Nombre completo: & \linead{10cm}\\
Fecha: & \linead{4cm} \quad Curso/grupo: \linead{4cm}\\
Mi firma: & \linead{9cm}\\
\end{tabularx}
\normalsize

\begin{infoband}{milcGradeDark}
\textbf{Para la próxima sesión.} Dos cosas. \textbf{Primero, haz tu remiendo}, el que escribiste con día y forma de verificarlo, y anota cómo te fue ---incluso si no lo hiciste; en ese caso anota qué se atravesó, que también es información---. \textbf{Segundo, pídele a alguien mayor que te muestre su revés:} busca a una persona que haga bien algo que a ti te parece difícil ---coser, cocinar, arreglar motores, manejar, tocar un instrumento, cortar el pelo--- y pregúntale \textbf{cómo le salían las primeras veces} y qué fue lo que más veces tuvo que rehacer. \emph{Vas a oír historias que no se cuentan: casi nadie va por ahí contando sus tres primeros intentos.} \textbf{Trae:} el resultado de tu remiendo y una frase de lo que te contaron.
\end{infoband}

\end{document}
